% numerals-ok: begin
\section{The three algorithms in pseudocode}
\label{app:algorithms}

The body says what each part does and why; this appendix says how it is run, at the level of a
reader who wants to re-implement it. Every constant named here is read from \texttt{data/*.json}
by the scripts in \texttt{code/}, and the three listings follow the control flow of
\texttt{p1a\_gpu\_segmentation.py}, \texttt{p2\_interval.py} and \texttt{p3\_closedform.py} in
that order.
% numerals-ok: end

\begin{algorithm}[H]
\caption{Boundary decoding: from per-frame scores to stage labels (Part I)}
\label{alg:p1}
\begin{algorithmic}[1]
\Require keypoint windows $X\in\mathbb{R}^{T\times 21\times 3}$, boundary network $f_\theta$,
smoothing $\sigma$, peak height $h$, minimum peak distance $d$
\Ensure stage label $y_t\in\{1,\dots,K\}$ for every frame
\State express each window in the head frame: subtract the head position, rotate by the
inverse head quaternion
\State $s_t \gets f_\theta(X)_t$ for $t=1,\dots,T$ \Comment{boundary score in $[0,1]$}
\State convolve $s$ with a Gaussian kernel of width $\sigma$
\State $P \gets \textsc{FindPeaks}(s;\;\text{height}=h,\;\text{distance}=d)$
\State $y_t \gets \#\{\,p\in P : p\le t\,\}$ \Comment{searchsorted, so labels stay contiguous}
\State \Return $y$, $P$
\end{algorithmic}
\end{algorithm}

The three decoding constants $(\sigma,h,d)$ are selected on validation records only and then
frozen; the alternative route in the body is chain-constrained Ward, which merges only adjacent
blocks and picks $K$ on the same validation records. The boundary F1 is computed by optimal
matching between predicted and annotated boundaries with a five-frame tolerance, so a decoder
that overshoots the count is penalised rather than rewarded.

\begin{algorithm}[H]
\caption{Ridge prediction and split-conformal radius (Part II)}
\label{alg:p2}
\begin{algorithmic}[1]
\Require context of $c=10$ observed frames, mask length $L\in\{6,15,30,60\}$ frames, training
windows and validation windows
\Ensure point prediction $\hat{y}\in\mathbb{R}^{3}$ and radius $q_L$
\State build the feature vector from the $c$ frames (positions, velocities, aperture, missing
flags) and standardise it with the training mean and variance
\State fit $\hat{\beta}_L=\arg\min_{\beta}\|Y-X\beta\|^{2}+\alpha\|\beta\|^{2}$ with
$\alpha=10$ \Comment{one model per mask length}
\State compute residuals $r_i=\|\hat{y}_i-y_i\|_{2}$ on the validation windows
\State $q_L\gets$ the $\lceil (n+1)(1-0.10)\rceil$-th smallest $r_i$
\Comment{split conformal, nominal coverage $90\%$}
\State \Return $\hat{y}$ and the ball $\{\,y:\|y-\hat{y}\|_2\le q_L\,\}$
\end{algorithmic}
\end{algorithm}

The radius is a single number per mask length rather than a per-axis interval, because the
error of a three-dimensional position is naturally measured by its norm. Coverage is evaluated
on test windows once, and the body reports both the radius and the realised coverage so that a
reader can see whether the nominal level was met.

\begin{algorithm}[H]
\caption{Reserve, closed-form punctuality, and dispatch (Part III)}
\label{alg:p3}
\begin{algorithmic}[1]
\Require case $c$, horizon $H$, budget $B$, review cap $K$, duration law $F_d$, and the
residual quantiles $Q_q$ of the duration model
\Ensure weighted on-time rate and cost
\State $t\gets 0$, $\mathit{cost}\gets 0$, $n_{\text{rev}}\gets 0$
\While{$t<H$ and dispatches $<200$}
  \State $\mathit{pressure}\gets \sum_{j\ \text{unstarted}}\hat{p}_j/(H-t)$
  \State $q\gets\operatorname{clip}(0.70-0.55\max(0,\mathit{pressure}-0.75),\,0.35,\,0.70)$
  \For{each feasible job $j$}
    \State $\hat{p}_j\gets\alpha\,\mathrm{pred}_j+\beta$; $D_j\gets$ travel distance;
    $s_j\gets t+D_j+0.5$
    \State $b_j\gets\max\!\big(0.2,\ \alpha(\mathrm{pred}_j+Q_q(r_j))+\beta\big)$
    \State $c_j\gets 2+0.1b_j+0.2D_j$; admit $j$ if $s_j+b_j\le H$ and $\mathit{cost}+c_j\le B$
  \EndFor
  \If{a second look is allowed and $n_{\text{rev}}<K$}
    \State compute $P_0$ and $P_1$ from the duration law at the current and next prefix
    \State request the look for the job maximising $w_j(P_1-P_0)-0.02$ if the maximum is positive
  \EndIf
  \State choose among admissible jobs the one maximising
  $w_j\max(0.03,P_j)/\max(0.5,\hat{p}_j)\cdot e^{-\mathrm{slack}_j/12}$,
  with $\mathrm{slack}_j=\mathrm{due}_j-s_j-\hat{p}_j$
  \State $\mathit{cost}\gets\mathit{cost}+c_j$; the job completes at $C_j=s_j+p_j$ with
  $p_j=\alpha d_j+\beta$; the crew is released at $s_j+\max(p_j,b_j)$
\EndWhile
\State \Return $\sum_j w_j\mathbf{1}[C_j\le\mathrm{due}_j]/\sum_j w_j$ and $\mathit{cost}$
\end{algorithmic}
\end{algorithm}

Line~5 is the only place where a forecast enters the decision, and it enters through the
quantile $Q_q$ of the residual law. The closed form of the body follows from line~11: since
$C_j=s_j+\alpha d_j+\beta$, the event $C_j\le\mathrm{due}_j$ is the event
$d_j\le(\mathrm{due}_j-s_j-\beta)/\alpha$, so its probability is one evaluation of $F_d$ at that
point and the reserve $b_j$ cancels out of it. That cancellation is why the paper can say that a
better forecast helps only through $F_d$, and why the paid second look, which changes
$s_j$ by one second without changing $F_d$, cannot buy punctuality.

\begin{table}[htbp]
\centering
\caption{The five columns a reader needs to audit one algorithm against the other two.}
\label{tab:algo-contract}
\small
\begin{tabular}{p{0.12\textwidth}p{0.18\textwidth}p{0.18\textwidth}p{0.18\textwidth}p{0.16\textwidth}}
\toprule
Part & Input & Output & Pass mark & Fails when \\
\midrule
I & keypoints of one record & stage labels, boundary set & ARI above the published decode & records shorter than one stage \\
II & ten observed frames & position and radius & coverage near $90\%$ and RMSE below hold-last & mask longer than the context carries information \\
III & duration law and job list & dispatch order and reserves & weighted on-time rate above the fixed-interval baseline & budget or horizon binds before the queue clears \\
\bottomrule
\end{tabular}
\end{table}
% numerals-ok: begin
Cost, in the sense the corpus reports it: Part~I is $O(T\log T)$ per record (one convolution
and one peak scan over $T$ frames); Part~II is $O(p^{3})$ once per mask length with $p$ the
feature count, and $O(p)$ per window at test time; Part~III is $O(J\log J)$ per dispatch event
with $J$ the queue length, and the grid search over reserves multiplies that by the number of
candidate quantiles. The measured wall-clock costs are in Table~\ref{tab:gpu-ledger} and
Appendix~G.
% numerals-ok: end

% numerals-ok: begin
The table is deliberately about failure conditions as well as pass marks: a method that cannot
say when it stops working is not yet a method. The body's Section~7 collects the same
conditions in prose, and Appendix~E records the attempts that failed them.
% numerals-ok: end
